\section{Задание 2}

\subsection{Методика измерения}

В программе \texttt{cache\_line.cpp} размер строки кеша оценивался по задержке зависимого обхода пар узлов. Строка --- блок данных, загружаемый в кеш при промахе. Если оба узла находятся в одной строке, при её загрузке вместе с первым узлом загружается второй. При переходе через границу строки второй узел занимает отдельную строку, поэтому число затронутых строк удваивается.

Каждая из 640 пар была размещена в отдельной области размером 384 байта и выравниванием 128 байт. Первый 8-байтный узел находился в начале области, второй --- по смещению $d$. Такое выравнивание размещает первый узел на границе строки при размере строки 64 байта. При обходе после первого узла посещался второй, затем выполнялся переход к первой позиции следующей пары. Для всех смещений использовалась одна случайная перестановка пар. Перед измерением проверялись порядок узлов и замыкание цикла (листинг~\ref{lst:cache-line-cycle}).

\verbatimlistingcaption{Размещение пар узлов и построение цикла}{lst:cache-line-cycle}
\begin{Verbatim}[frame=single,numbers=left,fontsize=\small,breaklines=true]
Алгоритм: построить_цикл_пар(d, порядок)
Вход: смещение d и порядок индексов 640 пар.
Выход: память и начальный адрес цикла из 1280 узлов.

если узлы перекрываются, выходят за границу области
или смещение нарушает выравнивание узла: ошибка
память <- 640 областей по 384 байта с выравниванием 128 байт
для каждой области с индексом i:
    первый[i] <- адрес узла в начале области
    второй[i] <- адрес узла по смещению d

для k от 0 до 639:
    i <- порядок[k]
    j <- порядок[(k + 1) по модулю 640]
    первый[i].следующий <- второй[i]
    второй[i].следующий <- первый[j]

начало <- первый[порядок[0]]
текущий <- начало
для k от 0 до 639:
    i <- порядок[k]
    если текущий не равен первый[i]: ошибка
    текущий <- текущий.следующий
    если текущий не равен второй[i]: ошибка
    текущий <- текущий.следующий
если текущий не равен началу: ошибка
вернуть память и начало
\end{Verbatim}

Для каждого смещения было проведено 10 испытаний по схеме задания 1 (листинги~\ref{lst:cache-levels-chase} и~\ref{lst:cache-levels-measure}): четыре прогревочных обхода и не менее 5~млн переходов в измеряемой серии. После каждого прогревочного обхода и измеряемой серии проверялся возврат к началу. За медиану было принято шестое значение отсортированного списка, то есть верхнее из двух центральных. Стандартное отклонение было вычислено с делением суммы квадратов отклонений на число испытаний, а коэффициент вариации --- как отношение отклонения к среднему в процентах (листинг~\ref{lst:cache-line-measure}).

\verbatimlistingcaption{Измерение задержки при фиксированном смещении}{lst:cache-line-measure}
\begin{Verbatim}[frame=single,numbers=left,fontsize=\small,breaklines=true]
Алгоритм: измерить_для_смещения(начало, n)
Вход: начальный адрес цикла и число узлов n.
Выход: статистики задержки в нс/шаг и вариация в процентах.

шаги <- n * округлить_вверх(5000000 / n)
задержки <- пустой список
повторить 10 раз:
    повторить 4 раза для прогрева:
        адрес <- обойти(начало, n)
        сохранить адрес и проверить возврат к началу
    t0 <- показание монотонных часов
    адрес <- обойти(начало, шаги)
    t1 <- показание монотонных часов
    сохранить адрес и проверить возврат к началу
    задержка <- (t1 - t0) в наносекундах / шаги
    добавить задержку в задержки
отсортировать задержки по возрастанию
медиана <- шестое значение
среднее <- среднее арифметическое задержек
отклонение <- стандартное отклонение задержек
вариация <- 100 * отклонение / среднее
вернуть медиану, среднее, отклонение и вариацию
\end{Verbatim}

Смещение изменялось от 8 до 128 байт с шагом 8 байт (листинг~\ref{lst:cache-line-main}).

\verbatimlistingcaption{Перебор смещений и обработка результатов}{lst:cache-line-main}
\begin{Verbatim}[frame=single,numbers=left,fontsize=\small,breaklines=true]
Алгоритм: эксперимент со смещениями внутри пары
Вход: заданные в коде смещения и число пар.
Выход: таблица статистик задержки для каждого смещения.

порядок <- перемешать индексы от 0 до 639
для d = 8, 16, ..., 128 байт:
    память, начало <- построить_цикл_пар(d, порядок)
    статистики <- измерить_для_смещения(начало, 1280)
    освободить память
    вывести d и статистики задержки
\end{Verbatim}

\subsection{Результаты измерений}

На рисунке~\ref{fig:cache-line} представлена зависимость медианной задержки от расстояния между узлами пары.

\begin{figure}[htbp]
    \centering
    \begin{tikzpicture}
        \begin{axis}[
            width=0.98\linewidth, height=7cm,
            xlabel={Расстояние между узлами пары, байт},
            ylabel={Задержка, нс/шаг},
            xtick={8,16,...,128},
            ytick={3.5,4,4.5,5,5.5},
            tick label style={font=\small},
            /pgf/number format/use comma,
            grid=major,
        ]
            \addplot[blue,thick,mark=*,mark size=1.5pt] coordinates {
                (8,3.647) (16,3.579) (24,3.617) (32,3.618)
                (40,3.556) (48,3.633) (56,3.430) (64,5.543)
                (72,5.422) (80,5.458) (88,5.403) (96,5.444)
                (104,5.540) (112,5.475) (120,5.443) (128,5.575)
            };
        \end{axis}
    \end{tikzpicture}
    \caption{Медианная задержка при разных расстояниях между узлами пары}
    \label{fig:cache-line}
\end{figure}

При увеличении смещения с 56 до 64 байт медианная задержка выросла с $3{,}430$ до $5{,}543$~нс/шаг. При больших смещениях она оставалась примерно на том же уровне. При строке размером 64 байта смещение 56 байт ещё позволяет целиком разместить второй 8-байтный узел в первой строке, тогда как при смещении 64 байта второй узел начинается со следующей строки. По этому переходу была получена оценка около 64 байт. В таблице~\ref{tab:cache-line} она сопоставлена со значением из столбца \texttt{COHERENCY-SIZE} вывода \texttt{lscpu -C} для L1d, L2 и L3.

\begin{table}[H]
    \centering
    \caption{Сравнение оценок размера строки кеша}
    \label{tab:cache-line}
    \begin{tabular}{|c|c|}
        \hline
        Оценка по графику & Размер по \texttt{lscpu} \\
        \hline
        $\approx 64$ байта & 64 байта \\
        \hline
    \end{tabular}
\end{table}

При 64-байтной строке обход затрагивает 40~КиБ кеш-строк, если узлы пары лежат в одной строке, и 80~КиБ, если в разных. Оба объёма больше L1d (32~КиБ), поэтому даже в первом случае все затронутые строки не помещаются в L1d.
